% Isolated post-CP36 candidate; no shared canonical or checkpoint edits.
\subsection{An actual Cartier tensor on the exceptional prime-five locus}

Work over an algebraically closed field of characteristic two.  Use
the prime-five tensor and matrices of (PHR1)--(PHR3).  In the chart
\(c_1=1,c_2=t\), let \(\sigma,\tau,E,Q\) denote the exact
polynomial vectors and quintics of (PEW1)--(PEW4).

\begin{lemma}[The exact prime-five common-wall criterion]
\label{lem:prime-five-common-wall-iff-ratio}
\status{Candidate; complete proof awaiting independent audits.}
For every nonzero prime-five Heisenberg tensor, the common
maximal-rank Euler wall is empty if and only if
\[
             c_1c_2\ne0,\qquad c_1^5=c_2^5.       \tag{PAE1}
\]
Here the common wall means the set of sections \([s]\) satisfying
\(\operatorname{rank}M_s=4\),
\(\operatorname{rank}B_{y_s}=4\), and
\(y_s(s^{[2]})=0\), where \(\ker M_s^t=ky_s\).
\end{lemma}
\begin{proof}
First suppose \(c_1\ne0\) and divide the tensor by this scalar.
This changes none of its kernels, ranks or wall equations, so we
may put \(c_1=1,c_2=t\).  By
Theorem~\ref{thm:prime-Heisenberg-birational-radical}, the five
maximal Pfaffians have greatest common divisor one.  Therefore
their common zero locus in \(\mathbf P^4_y\) has codimension at
least two.  Moreover
\[
        \frac{\partial Q}{\partial y_i}=\sigma_i(y)^2
                                                        \tag{PAE2}
\]
shows that the nonzero quintic \(Q\) is reduced: any repeated
irreducible factor of \(Q\) would divide every partial derivative,
contradicting the greatest-common-divisor assertion.

The exact polynomial certificate of
Lemma~\ref{lem:prime-five-exceptional-composition} proves that
\(E\) is a nonzero quintic for every \(t\), and that
\(E\) is proportional to \(Q\) exactly when \(t^5=1\).
If \(t^5\ne1\), then \(Q\nmid E\), because their degrees
agree.  As \(Q\) is reduced, it has an irreducible component
\(D\) on which \(E\) is not identically zero.  This component
cannot lie in the codimension-at-least-two common Pfaffian zero
locus.  Thus there is a point \([y]\in D\) with
\[
              E(y)\ne0,\qquad \sigma(y)\ne0.       \tag{PAE3}
\]
Put \(s=\sigma(y)\).  The Pfaffian condition gives
\(\operatorname{rank}B_y=4\), while
\(\tau(s)=E(y)y\ne0\) and
\(\operatorname{adj}(M_s)=s\tau(s)^t\) give
\(\operatorname{rank}M_s=4\).  Identity (PHR3) identifies its
left kernel with \(ky\).  Finally \(Q(y)=0\) is exactly
\(y(s^{[2]})=0\).  Hence the common wall is nonempty.
When \(t^5=1\), its emptiness is
Proposition~\ref{prop:prime-five-empty-common-wall}.

If \(c_1=0\), then \(c_2\ne0\).  Relabeling the theta basis
by \(e'_i=e_{2i}\), with the corresponding Frobenius basis in
the target, exchanges the two coefficient distances:
\(c'_1=c_2\), \(c'_2=c_{4}=c_1\).  The pairing
\(y(s^{[2]})\) is preserved by this simultaneous permutation.
After scalar normalization the new parameter is \(t=0\), for
which the preceding argument gives a nonempty wall.  The same
argument already treats \(c_2=0,c_1\ne0\).  This proves the
criterion in every coefficient chart.
\end{proof}

\begin{proposition}[An actual exceptional Cartier ratio]
\label{prop:prime-five-actual-exceptional-ratio}
\status{Candidate; complete proof awaiting independent audits.}
Let \(E\) be the supersingular elliptic curve
\(y^2+y=x^3\), put \(A=E\times E\), and write
\[
 V=\{O\}\times E,\quad H=E\times\{O\},\quad
 \Gamma=\{(P,-P):P\in E\},\qquad
                  L=\mathcal O_A(2V+H+\Gamma).     \tag{PAE4}
\]
In any compatible level-five theta basis in which its actual
nonzero Cartier tensor has form (PHR1), the two coefficients
satisfy (PAE1).  Consequently the exceptional empty-common-wall
locus in Proposition~\ref{prop:prime-five-empty-common-wall}
is realized by an actual ample globally generated
type-\((1,5)\) polarization on a nonordinary abelian surface.
\end{proposition}
\begin{proof}
Lemma~\ref{lem:root-elliptic-double-cover-lines} proves that
\(L\) is ample and globally generated of type \((1,5)\).
The actual tensor is nonzero by
Theorem~\ref{thm:DE-odd-polarization-Cartier-surjection} and has
form (PHR1) by Lemma~\ref{lem:prime-Heisenberg-two-contractions}.
By Theorem~\ref{thm:primitive-prime-generic-curve-rank}, the
smooth generic member has genus six and \(2\)-rank four.
Since \(A\) has \(2\)-rank zero, its relative Lang local
system is its full four-dimensional etale cohomology.

The geometric-generic double-cover calculation in
Theorem~\ref{thm:root-primitive-prime-permutation-counterexamples}
identifies this representation, naturally for Galois action,
with \(\mathbf F_2^{\mathcal B}/\mathbf F_2\), where the
branch set \(\mathcal B\) has five elements.  In particular
the geometric image is isomorphic to a subgroup of \(S_5\).
This assertion uses the actual generic curve and does not
require a finite etale branch scheme over an entire coefficient
open.

Suppose that the coefficient ratio did not satisfy (PAE1).
The preceding lemma produces a nonempty common maximal-rank
wall in the whole coefficient space \(X=\mathbf P H^0(A,L)\).
Lemma~\ref{lem:primitive-prime-Euler-criterion} identifies the
actual coherent relative Frobenius there with
\(M_s^t:(s^{[2]})^\perp\to s^\perp\), including sections
whose divisors are singular.  Thus all hypotheses of
Proposition~\ref{prop:birational-Cartier-kernel-wall-inertia}
hold on the common-rank open.  Its conclusion, with \(N=5\),
puts a subgroup
\[
          \operatorname{Hom}(\mathbf F_2,\mathbf F_2^3)
                        \simeq (\mathbf Z/2)^3       \tag{PAE5}
\]
in the generic geometric image.  The restriction to the dense
smooth sharp open does not change this generic image.  No
specialization of a full-rank Lang fibre across the wall is
being used.

For completeness \(S_5\) contains no elementary abelian
subgroup of order eight.  If an elementary abelian \(2\)-group
acts faithfully on five letters, its orbits have sizes powers
of two.  Its image on each orbit is regular, since an abelian
transitive action has its point stabilizer as its kernel.
The product of these orbit images has order at most four:
the only possibilities with a nontrivial orbit are partitions
\(4+1\), \(2+2+1\), or \(2+1+1+1\).  The faithful group
embeds in this product.  This contradicts (PAE5), so (PAE1)
must hold for the actual tensor.
\end{proof}

\begin{remark}[What the realization proves]
\label{rem:prime-five-actual-exceptional-ratio-scope}
\status{Candidate; actual realization only.}
The argument uses the exact polynomial identities already
certified in (PEW1)--(PEW4), the complete local moving-kernel
inertia proof, and the actual supersingular elliptic double-cover
calculation.  No new numerical check or unproved secant-variety
identification is used.  It does not compute an exact \(S_5\)
image, classify all type-\((1,5)\) Cartier tensors, or refute a
large-degree assertion for powers of a fixed embedding bundle.
It does show that the empty common wall is an actual geometric
phenomenon in this globally generated primitive family.
\end{remark}
